\clearpage\thispagestyle{otFront}\phantomsection\label{parcours}\pdfbookmark[0]{Construire son parcours}{bm-parcours}\otPageTitre{Construire son parcours}{Une année avec des retours réguliers}\fontsize{11}{15}\selectfont

\begin{tabularx}{\linewidth}{@{}lY@{}}\toprule
Nombres & 1 à 6 : milliards, décimaux, arrondis, produits, divisions, fractions quotients et calculs de fractions.\\[3mm]
Modéliser & 7 à 9 : pourcentages, schémas en barres, motifs, proportionnalité et échelles.\\[3mm]
Mesurer & 10 à 12 : longueurs et nombre $\pi$, aires et conversions, durées et heures décimales.\\[3mm]
Construire & 13 à 18 : distances, cercles, médiatrice, angles, triangles, symétrie axiale, solides et volumes.\\[3mm]
Organiser & 19 à 21 : enquêtes et données, probabilités, instructions et répétitions.\\\bottomrule
\end{tabularx}
\begin{savoirfaire}[Une séance qui fonctionne]Relire un exemple du cours,
chercher deux exercices sans corrigé, comparer, puis refaire la question
manquée. Vingt minutes régulières valent mieux qu’une longue séance
isolée.\end{savoirfaire}
\begin{savoirfaire}[Revenir et mélanger]Reprendre une notion quelques jours
plus tard, puis au moment des bilans. Les problèmes mélangent
volontairement les chapitres : c’est ainsi qu’on apprend à choisir.\end{savoirfaire}
Les défis demandent de 20 à 45 minutes, parfois en plusieurs fois.
Le tableau de bord suit les réussites et les reprises au fil de l’année.

\clearpage
